\section{Introduction}

Avant de développer RevMine, il est essentiel d'analyser les solutions existantes permettant la collecte et l'analyse des données de revue de code. Cette étude comparative nous permettra d'identifier les fonctionnalités disponibles, les lacunes à combler, et les meilleures pratiques à adopter.

Dans ce chapitre, nous examinerons différentes catégories d'outils de \textit{code review mining} : les outils de collecte de données, les plateformes d'analyse de métriques de processus, et les solutions académiques. Pour chaque outil, nous évaluerons ses fonctionnalités principales, ses points forts et ses limitations dans le contexte spécifique de l'extraction et de l'analyse des données de revue de code.

\section{Critères d'évaluation}

Pour structurer notre analyse, nous avons défini les critères d'évaluation suivants :

\begin{itemize}
    \item \textbf{Facilité d'utilisation :} L'outil est-il accessible aux non-experts ? Requiert-il des compétences techniques avancées ?
    \item \textbf{Couverture des plateformes :} Supporte-t-il GitHub, GitLab, ou d'autres plateformes de gestion de code ?
    \item \textbf{Automatisation :} La collecte et l'analyse sont-elles automatisées ou nécessitent-elles des interventions manuelles ?
    \item \textbf{Extensibilité :} Peut-on facilement ajouter de nouvelles analyses ou intégrer de nouvelles sources de données ?
    \item \textbf{Visualisation :} L'outil propose-t-il des tableaux de bord ou des visualisations interactives ?
    \item \textbf{Interactivité :} Permet-il aux utilisateurs de poser leurs propres questions ou d'explorer librement les données ?
    \item \textbf{Analyse de processus :} L'outil permet-il d'analyser le processus de revue (temps, interactions, patterns) ?
    \item \textbf{Maintenance et support :} L'outil est-il activement maintenu ? Dispose-t-il d'une communauté ou d'un support ?
\end{itemize}

\section{Outils de collecte de données}
Cette section présente les principaux outils dédiés à l’extraction et à la collecte de données issues des dépôts logiciels et des plateformes collaboratives. L’objectif est d’identifier leurs apports, leurs limites et leur adéquation avec les besoins d’une analyse des données de revue de code.

\subsection{PyDriller}

\textbf{PyDriller} \footnote{PyDriller Documentation officielle : \url{https://pydriller.readthedocs.io/en/latest/}}  est une bibliothèque Python populaire pour extraire des informations depuis des dépôts Git. Elle permet d'accéder aux commits, aux fichiers modifiés, aux auteurs, et aux métriques de complexité du code.

\textbf{Points forts :}
\begin{itemize}
    \item API Python simple et bien documentée
    \item Analyse locale des dépôts Git
    \item Extraction de métriques de code (complexité cyclomatique, lignes de code, etc.)
    \item Largement utilisé dans la recherche académique
    \item Support de l'analyse des commits et des modifications de fichiers
\end{itemize}

\textbf{Limitations :}
\begin{itemize}
    \item Ne collecte pas directement les données de revue de code (PR/MR, commentaires, discussions)
    \item Nécessite le clonage local du dépôt (problématique pour les grands projets)
    \item Pas d'interface graphique
    \item Aucun support natif pour les APIs GitHub/GitLab
    \item Limité à l'analyse Git, pas d'accès aux métadonnées de la plateforme
\end{itemize}

\subsection{Perceval}

\textbf{Perceval} \footnote{Perceval Dépôt officiel : \url{https://github.com/chaoss/grimoirelab-perceval}} (partie de GrimoireLab) est un outil Python spécialisé dans la collecte de données depuis des plateformes de développement collaboratif.

\textbf{Points forts :}
\begin{itemize}
    \item API Python claire et modulaire
    \item Support de multiples backends (GitHub, GitLab, Gerrit, Bugzilla, etc.)
    \item Collecte des pull requests, merge requests et commentaires
    \item Gestion automatique de l'authentification et de la pagination
    \item Bien documenté et activement maintenu
    \item Extraction incrémentale des données
\end{itemize}

\textbf{Limitations :}
\begin{itemize}
    \item Outil en ligne de commande (pas d'interface graphique)
    \item Nécessite des compétences en programmation Python
    \item Pas d'analyse intégrée (uniquement collecte)
    \item Configuration manuelle requise pour chaque source
    \item Sortie en format JSON brut nécessitant un post-traitement
\end{itemize}

\subsection{GHArchive}

\textbf{GHArchive} \footnote{GHArchive Site officiel : \url{https://www.gharchive.org/}} est un projet qui enregistre la chronologie publique de GitHub et fournit un accès aux événements historiques via Google BigQuery.

\textbf{Points forts :}
\begin{itemize}
    \item Accès à l'historique complet des événements GitHub depuis 2011
    \item Données structurées et prêtes à l'analyse
    \item Intégration avec BigQuery pour des requêtes SQL massives
    \item Idéal pour des études longitudinales à grande échelle
    \item Données actualisées quotidiennement
    \item Gratuit pour les requêtes (dans la limite du quota BigQuery)
\end{itemize}

\textbf{Limitations :}
\begin{itemize}
    \item Limité exclusivement à GitHub (pas de support GitLab ou autres)
    \item Nécessite des compétences en SQL et BigQuery
    \item Pas d'interface d'interrogation conviviale
    \item Coûts potentiels pour les analyses de grands volumes
    \item Données publiques uniquement (pas d'accès aux dépôts privés)
    \item Pas d'analyse ou visualisation intégrée
\end{itemize}

\subsection{GHTorrent}

\textbf{GHTorrent}\footnote{GHTorrent Page du projet : \url{https://ghtorrent.github.io/tutorial/}} est un projet académique qui crée une base de données miroir des données publiques de GitHub pour faciliter la recherche en génie logiciel.

\textbf{Points forts :}
\begin{itemize}
    \item Accès à des datasets historiques massifs
    \item Base de données relationnelle (MySQL) pour requêtes structurées
    \item Schéma de données bien documenté
    \item Utilisé dans de nombreuses publications académiques
    \item Dumps mensuels disponibles pour téléchargement
\end{itemize}

\textbf{Limitations :}
\begin{itemize}
    \item Limité à GitHub (pas de support GitLab)
    \item Projet en maintenance minimale depuis 2019
    \item Données historiques uniquement (pas de mises à jour continues)
    \item Nécessite des compétences en traitement de grandes bases de données
    \item Installation et configuration complexes
    \item Pas d'interface d'interrogation conviviale
\end{itemize}

\section{Plateformes d'analyse et de métriques}
Cette section s’intéresse aux plateformes qui vont au-delà de la simple collecte en proposant des mécanismes d’analyse, de calcul de métriques et de visualisation. Ces solutions permettent d’évaluer dans quelle mesure les données logicielles peuvent être transformées en indicateurs exploitables pour le suivi et l’amélioration des processus de développement.

\subsection{GrimoireLab}

\textbf{GrimoireLab} \footnote{GrimoireLab Site officiel : \url{https://github.com/chaoss/grimoirelab}} est une suite d'outils open source développée pour l'analyse complète de projets logiciels. Elle collecte des données depuis diverses sources (GitHub, GitLab, Jira, etc.) et les stocke dans une base de données Elasticsearch pour analyse et visualisation.

\textbf{Points forts :}
\begin{itemize}
    \item Supporte de nombreuses sources de données
    \item Collecte automatisée et incrémentale des PR/MR et commentaires
    \item Tableaux de bord Kibana pour la visualisation
    \item Métriques de processus de revue (temps de réponse, nombre de révisions)
    \item Utilisé par de grandes organisations (Mozilla, Eclipse Foundation)
    \item Architecture modulaire et extensible
\end{itemize}

\textbf{Limitations :}
\begin{itemize}
    \item Installation et configuration très complexes
    \item Nécessite une infrastructure complète (Elasticsearch, Kibana, MariaDB)
    \item Courbe d'apprentissage élevée
    \item Peu adapté aux analyses ad-hoc ou aux utilisateurs non techniques
    \item Pas d'interface conversationnelle ou de requêtes en langage naturel
    \item Consommation importante de ressources système
\end{itemize}

\subsection{Pluralsight Flow (ex-GitPrime)}

\textbf{Pluralsight Flow} \footnote{Flow Site officiel : \url{https://appfire.com/flow}} est une plateforme commerciale d'analyse de productivité des développeurs basée sur les données Git et les pull requests \cite{appfire_flow,appfire_flow_review_metrics}.

\textbf{Points forts :}
\begin{itemize}
    \item Métriques de productivité et d'efficacité des équipes
    \item Analyse du processus de revue (temps de cycle, temps de première revue)
    \item Visualisations sophistiquées et tableaux de bord
    \item Intégration avec GitHub, GitLab, Bitbucket
    \item Détection des goulots d'étranglement dans le processus de revue
    \item Interface web intuitive
\end{itemize}

\textbf{Limitations :}
\begin{itemize}
    \item Solution commerciale très coûteuse
    \item Boîte noire : métriques propriétaires non transparentes
    \item Pas d'accès aux données brutes
    \item Analyses prédéfinies uniquement, aucune personnalisation
    \item Préoccupations en matière de vie privée des développeurs
    \item Code source fermé, pas d'extensibilité
\end{itemize}

\subsection{Code Climate Velocity}

\textbf{Code Climate Velocity} \footnote{Code Climate Velocity Site officiel : \url{https://docs.velocity.codeclimate.com/}} est un module de la plateforme Code Climate dédié à l'analyse des métriques de développement et de revue de code.

\textbf{Points forts :}
\begin{itemize}
    \item Interface web intuitive
    \item Intégration avec GitHub et GitLab
    \item Métriques sur le processus de revue (taille des PR, temps de revue)
    \item Rapports de tendance et visualisations
    \item Corrélation entre qualité du code et processus de revue
\end{itemize}

\textbf{Limitations :}
\begin{itemize}
    \item Solution commerciale (coût élevé pour les grandes équipes)
    \item Analyses prédéfinies, peu de flexibilité
    \item Pas d'export facile des données brutes
    \item Fonctionnalités limitées comparées aux outils dédiés
    \item Pas de support pour requêtes personnalisées
\end{itemize}

\section{Solutions académiques et outils de recherche}

\subsection{RefDiff / RefactoringMiner } 

RefDiff\footnote{Dépôt officiel RefDiff : \url{https://github.com/aserg-ufmg/RefDiff}}  et RefactoringMiner\footnote{Dépôt officiel RefactoringMiner : \url{https://github.com/tsantalis/RefactoringMiner}} sont des outils académiques analysent les changements de code dans les commits et pull requests pour détecter des refactorings et comprendre l'évolution du code.

\textbf{Points forts :}
\begin{itemize}
    \item Détection précise de patterns de refactoring
    \item Utiles pour analyser la qualité des modifications dans les PR
    \item Support de plusieurs langages de programmation
    \item Utilisés dans de nombreuses études empiriques
\end{itemize}

\textbf{Limitations :}
\begin{itemize}
    \item Focus très spécifique (refactoring uniquement)
    \item Ne collectent pas les métadonnées de revue (commentaires, approbations)
    \item Outils de recherche avec maintenance irrégulière
    \item Documentation limitée
    \item Pas d'interface utilisateur
\end{itemize}

\subsection{SmartSHARK}

\textbf{SmartSHARK} \footnote{SmartSHARK Site officiel : \url{https://smartshark.github.io/}}  est une plateforme de recherche pour la collecte et l'analyse de données de projets logiciels, incluant les données de revue de code.

\textbf{Points forts :}
\begin{itemize}
    \item Architecture modulaire pour différentes sources de données
    \item Support de l'analyse de pull requests GitHub
    \item Base de données MongoDB pour stockage flexible
    \item Conçu pour la recherche reproductible
\end{itemize}

\textbf{Limitations :}
\begin{itemize}
    \item Projet académique avec maintenance limitée
    \item Installation et configuration complexes
    \item Documentation incomplète
    \item Communauté d'utilisateurs restreinte
    \item Pas d'interface graphique conviviale
\end{itemize}

\section{Tableau comparatif}

\begin{table}[H]
\centering
\small
\begin{tabularx}{\textwidth}{|l|Y|Y|Y|Y|Y|Y|}
\hline
\textbf{Outil} & \textbf{Collecte PR/MR} & \textbf{Analyse processus} & \textbf{Interface GUI} & \textbf{Langage naturel} & \textbf{Multi-plateforme} & \textbf{Open Source} \\
\hline
PyDriller & Non & Non & Non & Non & Partiel & Oui \\
\hline
Perceval & Oui & Non & Non & Non & Oui & Oui \\
\hline
GHArchive & Oui & Partiel & Non & Non & Non (GitHub) & Oui \\
\hline
GHTorrent & Oui & Partiel & Non & Non & Non (GitHub) & Oui \\
\hline
GrimoireLab & Oui & Oui & Oui (Kibana) & Non & Oui & Oui \\
\hline
Pluralsight Flow & Oui & Oui & Oui & Non & Oui & Non \\
\hline
Code Climate & Oui & Partiel & Oui & Non & Oui & Non \\
\hline
RefactoringMiner & Non & Non & Non & Non & N/A & Oui \\
\hline
SmartSHARK & Oui & Partiel & Non & Non & Partiel & Oui \\
\hline
\textbf{RevMine} & \textbf{Oui} & \textbf{Oui} & \textbf{Oui} & \textbf{Oui} & \textbf{Oui} & \textbf{Oui} \\
\hline
\end{tabularx}
\caption{Comparaison des outils de code review mining}
\label{tab:comparaison_outils}
\end{table}

\subsection{Comparaison détaillée : RevMine vs GrimoireLab}

GrimoireLab étant l'outil le plus complet identifié dans notre étude, nous proposons une comparaison approfondie avec RevMine pour mieux positionner notre contribution.

\subsubsection{Tableau comparatif}

\begin{tabularx}{\textwidth}{|l|X|X|}

\caption{Comparaison détaillée RevMine vs GrimoireLab}
\label{tab:revmine_vs_grimoirelab} \\

\hline
\textbf{Critère} & \textbf{GrimoireLab} & \textbf{RevMine} \\
\hline
\endfirsthead

\multicolumn{3}{c}%
{{\tablename\ \thetable{} -- Suite de la page précédente}} \\
\hline
\textbf{Critère} & \textbf{GrimoireLab} & \textbf{RevMine} \\
\hline
\endhead

\hline \multicolumn{3}{r}{\textit{Suite à la page suivante}} \\
\endfoot

\hline
\endlastfoot


\rowcolor{gray!25}
\multicolumn{3}{|l|}{\textbf{SOURCES DE DONNÉES}} \\
\hline
Plateformes supportées & 30+ (GitHub, GitLab, Gerrit, Bugzilla, Jira, Jenkins, etc.) & GitHub, GitLab (axé sur la revue de code) \\
\hline
Types de données & Multi-domaines (code, tickets, CI/CD, communication, documentation) & Spécialisé revue de code (PR/MR, commentaires, métriques processus) \\
\hline
Granularité métriques revue & Agrégées et générales & Fines et détaillées (15 catégories, 100+ métriques) \\
\hline

\rowcolor{gray!25}
\multicolumn{3}{|l|}{\textbf{MÉTRIQUES DE REVUE DE CODE}} \\
\hline
Métriques temporelles & Basiques (temps de fusion, dates création/fermeture) & Avancées (temps avant première revue, temps avant d'être prêt, temps en mode brouillon, HHO) \\
\hline
Métriques de participation & Contributeurs, réviseurs (identités unifiées) & Détaillées (auto-approuvé, auto-fusionné, engagement des réviseurs, contributeurs majeurs/mineurs) \\
\hline
Métriques de processus & Indice d'efficacité de revue, temps médians & Complètes (nombre révisions, itérations de commentaires, taille du remaniement, règles d'approbation) \\
\hline
Métriques de contenu & Non disponibles & Hashtags, @mentions, listes de tâches, tickets référencés \\
\hline
Métriques de conflit & Non disponibles & Possède des conflits, MR bloquantes, statut de fusion \\
\hline
Métriques LLM-dérivées & Non disponibles & Classification types revue, détection de refactorisation, indicateurs de qualité de code \\
\hline

\rowcolor{gray!25}
\multicolumn{3}{|l|}{\textbf{INTERFACE \& UTILISABILITÉ}} \\
\hline
Interface utilisateur & Kibana (tableaux de bord pré-configurés) & React (interface guidée + conversationnelle) \\
\hline
Courbe apprentissage & Très élevée (Elasticsearch, Kibana, configuration complexe) & Modérée (interface intuitive, guidage LLM) \\
\hline
Requêtes personnalisées & Langage de requête Kibana (expertise requise) & Langage naturel via LLM \\
\hline
Configuration collecte & Fichiers YAML/JSON complexes & Interface graphique + génération par LLM \\
\hline
Public cible & Experts DevOps, administrateurs système & Chercheurs, praticiens, managers (tous niveaux) \\
\hline

\rowcolor{gray!25}
\multicolumn{3}{|l|}{\textbf{ANALYSE \& EXPLORATION}} \\
\hline
Module LLM & Non & Oui (génération plans, analyses, requêtes LN) \\
\hline
Analyses prédéfinies & Oui (panels Sigils) & Oui + personnalisables via LLM \\
\hline
Nettoyage données & Minimal & Moyen \\
\hline
Export données & Exports Elasticsearch, JSON & CSV, JSON, SQL + reproductibilité garantie \\
\hline

\rowcolor{gray!25}
\multicolumn{3}{|l|}{\textbf{CAS D'USAGE}} \\
\hline
Usage principal & Monitoring continu multi-projets, tableaux de bord organisationnels & Études empiriques ciblées, analyses ad-hoc, optimisation processus revue \\
\hline
Adapté recherche & Oui (mais expertise technique requise) & Oui (axé sur reproductibilité, accessibilité) \\
\hline
Adapté industrie & Oui (déploiements production, ex: Wikimedia) & Oui (analyse équipes, identification goulots d'étranglement) \\

\end{tabularx}
\subsubsection{Analyse comparative}

Cette comparaison met en évidence deux philosophies distinctes. GrimoireLab s'impose comme une solution industrielle robuste, conçue pour un monitoring organisationnel à très grande échelle et capable de traiter plus de 30 types de sources de données (code, CI/CD, communication, etc.) \cite{duenas2021grimoirelab}. Cette puissance, documentée dans l'étude, s'accompagne toutefois d'une complexité de déploiement et d'utilisation élevée (Elasticsearch, Kibana), et ses métriques de revue de code restent générales \cite{duenas2021grimoirelab}. À l'inverse, RevMine se positionne comme un outil spécialiste, focalisé exclusivement sur la revue de code (GitHub/GitLab). Il sacrifie la largeur des sources pour offrir une granularité d'analyse supérieure (plus de 100 métriques) et, surtout, une accessibilité radicalement améliorée grâce à une interface graphique dédiée et l'intégration d'un LLM pour les requêtes en langage naturel. RevMine comble ainsi une lacune en matière d'analyse sémantique approfondie et d'accessibilité pour les praticiens et les chercheurs, là où GrimoireLab excelle dans le monitoring de haut niveau.

\section{Synthèse et positionnement de RevMine par rapport au outils}

L'analyse des solutions existantes révèle plusieurs constats importants :

\begin{itemize}
    \item  \textbf{Fragmentation des solutions} \\
    Les outils existants se concentrent généralement sur un aspect spécifique du code review mining : soit la collecte de données brutes (PyDriller, Perceval, GHArchive), soit l'analyse de métriques prédéfinies (GrimoireLab, Pluralsight Flow). Aucune solution n'offre une approche intégrée couvrant l'ensemble du pipeline depuis la collecte jusqu'à l'analyse interactive personnalisée.
    \\
    \item  \textbf{Barrière technique élevée} \\
    La plupart des outils open source (PyDriller, Perceval, GrimoireLab, GHArchive) nécessitent des compétences techniques avancées en programmation, en administration système ou en requêtes SQL/NoSQL. Les outils commerciaux sont plus accessibles mais offrent peu de flexibilité, sont coûteux et ne permettent pas d'accéder aux données brutes.
    \\
    \item  \textbf{Manque d’interactivité et de flexibilité} \\
    Les plateformes existantes proposent généralement des analyses et des visualisations prédéfinies sans permettre aux utilisateurs de formuler leurs propres questions ou d’explorer librement les données selon leurs besoins spécifiques de recherche ou d’analyse.
    \\
    \item  \textbf{bsence d’interface conversationnelle} \\
    Aucun outil de code review mining n’intègre de capacité d’interrogation en langage naturel. Les utilisateurs doivent soit écrire du code (Python, SQL), soit se limiter
    aux tableaux de bord prédéfinis, ce qui limite considérablement l’accessibilité et la flexibilité d’analyse.
    \\
    \item  \textbf{Limites de scalabilité ou d’accessibilité} \\
    Les outils académiques (GHTorrent, SmartSHARK) offrent un accès à de grands volumes de données mais sont difficiles à installer et maintenir. Les solutions commerciales sont plus accessibles mais limitées en termes de transparence et d’extensibilité.
    
\end{itemize}

\subsection{Positionnement}

Face à ces lacunes, \textbf{RevMine} se positionne comme une solution innovante combinant :

\begin{enumerate}
    \item \textbf{Collecte automatisée complète} : Module robuste pour extraire l'ensemble des données de revue de code (PR/MR, commentaires, révisions, approbations) depuis GitHub et GitLab avec gestion complète des contraintes API (rate limiting, pagination, authentification)
    
    \item \textbf{Interface accessible} : GUI intuitive permettant aux non-experts de configurer et lancer des collectes sans programmation, tout en offrant des options avancées pour les utilisateurs expérimentés
    
    \item \textbf{Analyse de processus} : Calcul automatique de métriques de processus de revue (temps de réponse, nombre de révisions, patterns de collaboration) avec possibilité d'ajouter de nouvelles métriques
    
    \item \textbf{Exploration interactive} : Tableaux de bord personnalisables pour l'exploration visuelle des données collectées
    
    \item \textbf{Requêtes en langage naturel} : Intégration d'un LLM permettant aux utilisateurs de poser des questions dans leur langue naturelle et d'obtenir des analyses contextualisées sans avoir à écrire de code ou de requêtes SQL
    
    \item \textbf{Approche open source} : Code source disponible, architecture modulaire extensible et bien documentée pour encourager la contribution communautaire et la reproductibilité de la recherche
    
    \item \textbf{Focus recherche et pratique} : Conçu pour répondre aux besoins des chercheurs en ingénierie logicielle empirique (extraction de données pour publications, analyses reproductibles) tout en restant accessible aux praticiens (chefs de projet, responsables techniques)
\end{enumerate}

\section{Conclusion}

L'analyse des solutions existantes confirme l'absence d'une solution complète, accessible et interactive pour le code review mining. Les outils existants présentent des fonctionnalités partielles mais aucun ne combine l'ensemble des caractéristiques nécessaires pour démocratiser l'accès à ces analyses :

\begin{itemize}
    \item Les outils de collecte (Perceval, GHArchive) sont puissants mais nécessitent des compétences techniques et ne fournissent pas d'analyse intégrée
    \item Les plateformes d'analyse (GrimoireLab, Pluralsight Flow) offrent des visualisations mais sont soit complexes à déployer, soit limitées aux analyses prédéfinies
    \item Les solutions académiques sont trop spécialisées ou insuffisamment maintenues
    \item Aucune solution n'offre d'interface conversationnelle pour l'interrogation en langage naturel
\end{itemize}

RevMine vise à combler ces lacunes en proposant une approche intégrée qui tire parti des meilleures pratiques identifiées (robustesse de Perceval pour la collecte, visualisation de GrimoireLab) tout en innovant avec l'intégration d'un module LLM pour l'interrogation en langage naturel, le tout dans une interface accessible.

Le chapitre suivant détaillera l'analyse des besoins, la conception architecturale et la réalisation concrète de RevMine.